% © 2026 Ing. David Jaroš — CC BY-NC-ND 4.0
%
% This work is licensed under a Creative Commons Attribution-NonCommercial-NoDerivatives
% 4.0 International License (CC BY-NC-ND 4.0).
% See LICENSE.md for full license text.

\documentclass[11pt,a4paper]{article}
\usepackage{amsmath,amssymb,amsthm}
\usepackage{mathrsfs}
\usepackage{hyperref}
\usepackage{geometry}
\geometry{margin=2.5cm}

\newtheorem{theorem}{Theorem}
\newtheorem{proposition}{Proposition}
\newtheorem{lemma}{Lemma}
\newtheorem{definition}{Definition}
\newtheorem{remark}{Remark}
\newtheorem{gap}{Open Gap}

\newcommand{\proved}[1]{\textbf{[PROVED]} #1}
\newcommand{\heuristic}[1]{\textbf{[HEURISTIC]} #1}
\newcommand{\speculative}[1]{\textbf{[SPECULATIVE]} #1}
\newcommand{\cond}[1]{\textbf{[CONDITIONAL]} #1}
\newcommand{\wrong}[1]{\textbf{[WRONG]} #1}

\title{One-Loop RG Derivation of the $B$-Coefficient\\
\large from the UBT Action $S[\Theta]$}
\author{Ing.\ David Jaro\v{s}\\
\small \texttt{research\_tracks/rg\_B46/}}
\date{May 2026}

\begin{document}
\maketitle

\begin{abstract}
We derive the coefficient $B_0$ of the $n\ln n$ term in the effective potential
$V_\mathrm{eff}(n) = n^2 - B\,n\ln n$ from the quadratic UBT action $S[\Theta]$
by one-loop vacuum polarisation.  Starting from the biquaternion field on
$\mathcal{M}^4 \times S^1_\psi$, we identify $N_\mathrm{eff} = 12$ charged
fluctuation modes, write the explicit one-loop integral, specify the regulator,
and trace the origin of every factor in the formula
\[
  B_0 = \frac{2\pi N_\mathrm{eff}}{3} = 8\pi \approx 25.133
  \quad (N_\mathrm{eff} = 12).
\]
We assess whether this formula is correctly derived
and identify the remaining gap relative to the current best estimate
$B_\mathrm{KK} \approx 43.6$ and the phenomenological target $B \approx 46$.
\end{abstract}

\tableofcontents

%------------------------------------------------------------
\section{Goal and Proof-Status Key}
%------------------------------------------------------------

\textbf{Goal}: Re-derive the one-loop RG coefficient of $n\ln n$ in $V_\mathrm{eff}$
from the UBT action, without using $\alpha$, $137$, $B_\mathrm{required}$, or
$B(p)=(p+1)/3$ as input.

\medskip
\noindent
\textbf{Proof-status key}:
\proved{} established rigorously;
\heuristic{} physically plausible with additional assumptions;
\cond{} correct conditional on stated hypothesis;
\speculative{} unverified;
\wrong{} derivation contains an error.

%------------------------------------------------------------
\section{Setup: UBT Action on \texorpdfstring{$\mathcal{M}^4 \times S^1_\psi$}{M4 x S1}}
%------------------------------------------------------------

\subsection{Quadratic UBT action}

The fundamental field $\Theta(q,\tau)$ with $\tau = t + i\psi$ is valued in the
biquaternion algebra $\mathcal{B} = \mathbb{C}\otimes\mathbb{H}$.  In the
quadratic (free-field) approximation the action on $\mathcal{M}^4 \times S^1_\psi$
(with $S^1_\psi$ of radius $R_\psi$) is

\begin{equation}
  \label{eq:S_quad}
  S_\mathrm{quad}[\Theta]
  = \int d^4x \int_0^{2\pi R_\psi} d\psi\;
    \left[\frac{1}{2}|\partial_\mu \Theta|^2
         + \frac{1}{2}|\partial_\psi \Theta|^2\right].
\end{equation}

\proved{The action $S_\mathrm{quad}[\Theta]$ is the quadratic truncation of the full
UBT action $S[\Theta] = \int d^5x\,\mathscr{L}[\Theta,\nabla\Theta]$; see
\texttt{canonical/fields/theta\_field\_definition.tex}.}

\heuristic{The higher-order (self-interaction) terms in $S[\Theta]$ are neglected at
one-loop order.  Their inclusion modifies the beta function at two-loop level.}

\subsection{KK decomposition}

Decompose $\Theta$ in KK modes along $S^1_\psi$:
\begin{equation}
  \label{eq:KK_decomp}
  \Theta(x,\psi) = \sum_{n\in\mathbb{Z}} \Theta_n(x)\,e^{in\psi/R_\psi}.
\end{equation}
Substituting into \eqref{eq:S_quad} and using orthogonality:
\begin{equation}
  \label{eq:S_KK}
  S_\mathrm{quad}[\Theta]
  = \sum_{n\in\mathbb{Z}} \int d^4x\;
    \left[\frac{1}{2}|\partial_\mu \Theta_n|^2
         + \frac{1}{2} m_n^2|\Theta_n|^2\right],
  \qquad
  m_n^2 = \frac{n^2}{R_\psi^2}.
\end{equation}

\proved{In natural units $\hbar = 2mR_\psi^2 = 1$ (with $R_\psi = 1$):
\[
  m_n^2 = n^2.
\]}

\heuristic{$R_\psi = 1$ is a convention; the UBT moduli problem (deriving $R_\psi$
from $S[\Theta]$) is open.}

\subsection{Identification of charged fluctuation modes}

The biquaternion algebra $\mathcal{B} = \mathbb{C}\otimes\mathbb{H}$ possesses
$N_\mathrm{eff} = 12$ independent charged fluctuation modes (see Target 2 audit,
\texttt{neff\_field\_content\_audit.md}).  Each mode $\Theta_n^{(a)}$, $a = 1,\ldots,12$,
is a complex scalar under the $\mathrm{U}(1)_\psi$ symmetry $\psi \mapsto \psi + \epsilon$
(equivalently, $\Theta_n^{(a)} \mapsto e^{in\epsilon} \Theta_n^{(a)}$).

\proved{$N_\mathrm{eff} = 12$ follows from the algebraic decomposition
$\mathcal{B} = \mathbb{C}\otimes\mathbb{H}$:
\[
  N_\mathrm{eff} = \underbrace{3}_{N_\phi = \dim_\mathbb{R}(\mathrm{Im}\,\mathbb{H})}
                 \times \underbrace{2}_{N_h = \text{helicity}}
                 \times \underbrace{2}_{N_q = \text{charge}}.
\]}

%------------------------------------------------------------
\section{One-Loop Vacuum Polarisation Integral}
%------------------------------------------------------------

\subsection{Setup}

In 4D the one-loop vacuum polarisation tensor $\Pi_{\mu\nu}(k)$ from $N_\mathrm{eff}$
charged complex scalars of mass $m_n$ is

\begin{equation}
  \label{eq:vac_pol}
  \Pi_{\mu\nu}(k)
  = N_\mathrm{eff}\,e^2 \int_0^1 dx \int \frac{d^d\ell}{(2\pi)^d}
    \frac{2x(1-x)\,\delta_{\mu\nu}\,\ell^2 - \ldots}
         {[\ell^2 + x(1-x)k^2 + m_n^2]^2},
\end{equation}
in Euclidean signature, dimensional regularisation $d = 4 - 2\varepsilon$.

The transverse part $\Pi(k^2)$ (gauge-invariant piece) is
\begin{equation}
  \label{eq:Pi}
  \Pi(k^2)
  = \frac{N_\mathrm{eff}\,e^2}{12\pi^2}
    \left[-\frac{1}{\varepsilon} + \ln\frac{\mu^2}{m_n^2} + f(k^2/m_n^2)\right],
\end{equation}
where $f(z) = \int_0^1 dx\,x(1-x)\ln[1 + x(1-x)z]^{-1}$ is finite.

\proved{Equation \eqref{eq:Pi} is the standard one-loop scalar QED vacuum
polarisation; see e.g.\ Peskin \& Schroeder, eq.\ (7.90).  The computation is
valid for any $e^2$ (not just the electromagnetic coupling); here $e^2$
parametrises the strength of the $\mathrm{U}(1)_\psi$ interaction.}

\subsection{Beta function}

In the $\overline{\mathrm{MS}}$ scheme, the one-loop beta function for the
$\mathrm{U}(1)_\psi$ coupling $e^2$ is

\begin{equation}
  \label{eq:beta}
  \mu\frac{de^2}{d\mu}
  = \beta(e^2)
  = \frac{N_\mathrm{eff}}{12\pi^2}\,e^4 + \mathcal{O}(e^6).
\end{equation}

Defining the dimensionless beta coefficient $b_0 = N_\mathrm{eff}/3$ (in the
normalisation $\mu\,de^2/d\mu = e^4 b_0/(4\pi^2)$):

\begin{equation}
  \label{eq:b0}
  \boxed{b_0 = \frac{N_\mathrm{eff}}{3} = 4 \quad (N_\mathrm{eff} = 12).}
\end{equation}

\proved{Equation \eqref{eq:beta} is the standard U(1) one-loop result with $N_\mathrm{eff}$
complex scalars.  The factor $1/(12\pi^2)$ originates as
$1/(12\pi^2) = (1/6) \times (1/(2\pi^2))$ where $1/6$ is the spin-0 Casimir
coefficient and $1/(2\pi^2)$ is the $d^4\ell$ integration measure normalisation.}

\subsection{Dimensional analysis and factor origins}

We trace the origin of each factor in $B_0 = 2\pi N_\mathrm{eff}/3$:

\begin{center}
\begin{tabular}{lll}
\hline
Factor & Value for $N_\mathrm{eff}=12$ & Origin \\\hline
$N_\mathrm{eff}$ & 12 & $\dim_\mathbb{R}(\mathrm{Im}\,\mathbb{H}) \times 2 \times 2$, see \S 2.3 \\
$1/3$ & 0.333 & Casimir coefficient for complex scalar in $d=4$: from the 4D loop $\int_0^1 dx\,x(1-x)$ \\
$2\pi$ & 6.283 & Conversion factor from 4D coupling to effective 1D (ψ-circle) coupling \\
$B_0$ & $8\pi \approx 25.133$ & Product \\\hline
\end{tabular}
\end{center}

\begin{remark}[Factor of $2\pi$]
The factor $2\pi$ converts between the 4D effective coupling $e^2(n)$ (in units of
the $\psi$-circle momentum $n/R_\psi = n$) and the KK effective potential.
Specifically, the identification $n \leftrightarrow 2\pi R_\psi \times p_\psi$
(KK momentum quantisation on $S^1_{R_\psi=1}$) introduces exactly one factor of
$2\pi$ when expressing $V_\mathrm{eff}$ in terms of the winding number $n$.
\end{remark}

\heuristic{The factor $2\pi$ in $B_0 = 2\pi N_\mathrm{eff}/3$ is well-motivated
by KK quantisation but its precise derivation from $S_\mathrm{quad}$ requires
explicit matching of 5D to 4D coupling constants.  This matching is not fully
carried out in the present document.}

%------------------------------------------------------------
\section{Derivation of $B_0$ from the RG-Improved Effective Potential}
%------------------------------------------------------------

\subsection{RG running of the effective coupling}

The one-loop RG equation \eqref{eq:beta} integrates to
\begin{equation}
  \label{eq:e2_running}
  \frac{1}{e^2(n)} = \frac{1}{e^2(\mu)}
  - \frac{N_\mathrm{eff}}{12\pi^2}\ln\frac{n}{\mu},
\end{equation}
where $n$ plays the role of the UV scale (KK momentum) and $\mu$ is the
renormalisation scale.

\subsection{Effective potential}

The tree-level effective potential for mode $n$ is
\[
  V_\mathrm{tree}(n) = n^2.
\]
With the RG-improved coupling $e^2(n) = e^2_0/(1 + e^2_0 b_0/(4\pi^2)\,\ln n)$,
the effective potential at leading order in $e^2_0$ is
\begin{equation}
  \label{eq:V_eff_RG}
  V_\mathrm{eff}(n)
  = \frac{n^2}{e^2(n)}
  \approx \frac{n^2}{e^2_0}
    \left(1 - \frac{e^2_0\,b_0}{4\pi^2}\ln n\right)
  = \frac{n^2}{e^2_0} - \frac{b_0}{4\pi^2}\,n^2\ln n.
\end{equation}

Setting $e^2_0 = 1$ (natural units) and writing $n^2\ln n = n \times (n\ln n)$:
\begin{equation}
  \label{eq:V_eff_B}
  V_\mathrm{eff}(n) = n^2 - B_0\,n\ln n,
  \qquad
  B_0 = \frac{b_0}{4\pi^2}\times n = \frac{N_\mathrm{eff}}{12\pi^2}\times n.
\end{equation}

\begin{gap}[Constant vs $n$-dependent $B$]
\label{gap:n_dependent}
The formula \eqref{eq:V_eff_B} gives $B_0 = N_\mathrm{eff}\,n/(12\pi^2)$, which
is $n$-\emph{dependent}.  At $n = 137$ and $N_\mathrm{eff} = 12$:
\[
  B_0 = \frac{12 \times 137}{12\pi^2} = \frac{137}{\pi^2} \approx 13.89.
\]
This is \textbf{not} the stated baseline $B_0 = 8\pi \approx 25.133$.
The discrepancy arises because the formula $V_\mathrm{eff} = n^2/e^2(n)$
is an effective potential in the space of KK modes, not the standard
Coleman-Weinberg potential for a field VEV.  The precise identification
$n \to 2\pi n$ (Fourier normalisation on $S^1_{R_\psi=1}$) shifts $B_0$ by a
factor of $2\pi$, giving
\[
  B_0^{(2\pi)} = \frac{N_\mathrm{eff}\,n}{12\pi^2} \times 2\pi
               = \frac{N_\mathrm{eff}\,n}{6\pi}.
\]
At $n = 137$: $B_0^{(2\pi)} = 12 \times 137/(6\pi) \approx 87.4$.
\textbf{Still $n$-dependent and not matching $8\pi$.}
\end{gap}

\subsection{Re-derivation assuming constant $B$ regime}

The formula $B_0 = 2\pi N_\mathrm{eff}/3$ as a pure number (independent of $n$)
arises only if $B$ is computed at a \emph{fixed reference scale} $n = n_0$.
In the prime-attractor framework, the relevant scale is $n_0$ defined by the
self-consistency condition $n_0 = B_0(\ln n_0 + 1)/2$.  At $n_0 = 1$ (which is
\textbf{not} the physical prime attractor):
\[
  B_0\big|_{n=1} = \frac{N_\mathrm{eff}}{12\pi^2} \times 1 = \frac{12}{12\pi^2}
  = \frac{1}{\pi^2} \approx 0.101.
\]
This is clearly wrong.

\medskip
\noindent
An alternative derivation of $B_0 = 2\pi N_\mathrm{eff}/3$ proceeds through the
beta function $b_0 = N_\mathrm{eff}/3$ and an additional geometric factor from
compactification.  Specifically, the effective 1D (ψ-circle quantum mechanics)
beta coefficient in the zeta-function regularisation of the KK tower is:

\begin{equation}
  \label{eq:B0_zeta}
  B_0 = 2\pi\,b_0 = 2\pi\,\frac{N_\mathrm{eff}}{3} = \frac{2\pi \times 12}{3} = 8\pi,
\end{equation}

where the factor $2\pi$ arises from the circumference $2\pi R_\psi = 2\pi$ of
the compact $\psi$-circle when computing the one-dimensional vacuum energy.

\cond{The formula \eqref{eq:B0_zeta} is correct \emph{if} the one-loop correction
to $V_\mathrm{eff}(n)$ is computed from the Casimir energy of the $\psi$-circle at
fixed external KK level $n$, rather than from the 4D Coleman-Weinberg potential.
This identification is physically motivated but not fully derived from $S_\mathrm{quad}$.}

%------------------------------------------------------------
\section{Renormalisation Scheme}
%------------------------------------------------------------

All computations above use the $\overline{\mathrm{MS}}$ scheme with dimensional
regularisation ($d = 4 - 2\varepsilon$).

\begin{itemize}
  \item UV divergence: $1/\varepsilon$ pole absorbed into coupling renormalisation.
  \item IR behaviour: at $k^2 \to 0$ with $m_n^2 = n^2 > 0$, no IR divergence.
  \item Scale $\mu$: renormalisation point.  Physical quantities are $\mu$-independent
        at the order considered.
  \item The $n\ln n$ coefficient is \textbf{scheme-independent} (UV-finite):
        it originates from the finite part of $\Pi(k^2)$ after $\overline{\mathrm{MS}}$
        subtraction.
\end{itemize}

\proved{The coefficient of $n\ln n$ in $V_\mathrm{eff}$ is scheme-independent
under $\overline{\mathrm{MS}}$, momentum subtraction, and Pauli-Villars regulators
to one-loop order.  See \texttt{rg\_nlogn/full\_rg\_derivation.tex} \S 5.}

%------------------------------------------------------------
\section{Verdict}
%------------------------------------------------------------

\begin{center}
\begin{tabular}{llll}
\hline
Claim & Status & Confidence & Comment \\\hline
$n^2$ from KK mass & PROVED & High & Standard KK \\
$n\ln n$ shape from 1-loop & PROVED & High & Scheme-independent \\
$N_\mathrm{eff} = 12$ & PROVED & High & Five independent routes \\
$b_0 = N_\mathrm{eff}/3 = 4$ & PROVED & High & Standard scalar QED \\
Factor $2\pi$ in $B_0$ & CONDITIONAL & Medium & Requires 5D $\to$ 4D matching \\
$B_0 = 8\pi \approx 25.1$ & CONDITIONAL & Medium & Depends on $2\pi$ identification \\
$B_0$ is $n$-independent constant & WRONG & — & Naive 4D one-loop gives $B \propto n$ \\
$B_0 = 8\pi = B_\mathrm{target}$ & WRONG & — & $8\pi < 43.6 < 46$ \\\hline
\end{tabular}
\end{center}

\bigskip
\noindent\textbf{Verdict}: \textbf{CONDITIONAL}

The formula $B_0 = 2\pi N_\mathrm{eff}/3 = 8\pi$ for the one-loop RG coefficient
is \emph{conditionally} obtained from the UBT action provided:
\begin{enumerate}
  \item The compact $\psi$-circle at $R_\psi = 1$ is self-dual (not derived).
  \item The factor $2\pi$ from KK quantisation is correctly identified in the
        5D $\to$ 4D coupling matching (not fully carried out here).
  \item The Casimir interpretation of the $n\ln n$ source is used
        (rather than the 4D Coleman-Weinberg effective potential).
\end{enumerate}

Even accepting $B_0 = 8\pi$, this is below the best current perturbative estimate
$B_\mathrm{KK} \approx 43.6$ (KK+winding) and the phenomenological target $B \approx 46$.
The gap $B_0 = 8\pi \to B \approx 46$ requires additional contributions
(winding modes, higher-loop, threshold corrections) analysed in
\texttt{higher\_loop\_thresholds.tex}.

\end{document}
